\section{Introduction}
\label{sec:intro}
Post-training with reinforcement learning (RL) is what turns a pretrained autoregressive (AR) language model into a model that reasons \citep{ouyang2022instructgpt,shao2024deepseekmath,deepseek2025r1}. The language model is a policy whose action is the next token: at every position it outputs a categorical distribution over the vocabulary and samples one token from it. A completion is sampled, a reward is obtained from a verifier or a reward model, an advantage is formed against a baseline or a group of samples, and the score-function policy gradient \citep{williams1992reinforce} raises the log-probability of every token of a rewarded completion and lowers it for an unrewarded one, with PPO or GRPO supplying the trust region and the baseline \citep{schulman2017ppo,shao2024deepseekmath}. This works because the action space is made for it: every action is a discrete, low-dimensional, semantically meaningful decision with an explicit probability, and credit attaches to exactly those decisions.

A new family of language models takes no token actions at all. Continuous flow language models such as ELF \citep{hu2026elf} represent a sequence as real-valued token embeddings and generate it by flow matching \citep{lipman2022flow,liu2022rectified}: starting from Gaussian noise, a denoiser transports the whole sequence toward clean embeddings over a number of \emph{flow steps}, and tokens are read out only at the end. Sampling is a chain of Gaussian transitions over an $L\times d$ real-valued state, there is no token distribution to sample from along the way, and the entire chain is differentiable. How to post-train such a model is open. The tempting answer is to port the AR recipe by treating each denoising step as an action with a Gaussian likelihood, as has been done for image flows \citep{black2023ddpo,liu2025flowgrpo,xue2025dancegrpo} and for masked diffusion language models \citep{zhao2025d1,zhu2025llada15}.

We study this question at three levels. The first is the port itself. The score-function identity applies verbatim to the Gaussian chain, and GRPO does improve the flow model on every task we try. But the ``actions'' it differentiates are not decisions; they are the noise coordinates injected at every step, $L\times d$ of them per step and $K$ steps per sample, none of which decides a token, because a token is committed only by the accumulated dynamics of the chain. The score of a Gaussian step is the injected noise times a Jacobian, so every reward bit is averaged over an astronomically larger set of nuisance directions than in the AR case, where it attaches to a few dozen categorical decisions (\S\ref{sec:estimators}). No refinement of the advantage changes this, and it shows in the results as gains that are smallest where the headroom is largest and that vary across seeds.

The second level uses what the flow model has and the AR model lacks: differentiability. The pathwise (reparameterised) gradient is absent from AR post-training only because categorical sampling cannot be reparameterised; here it is the native estimator, the one used in continuous control and in reward back-propagation for image diffusion \citep{lillicrap2015ddpg,haarnoja2018sac,prabhudesai2023alignprop,clark2023draft}. Because the reward exists only on the decoded text, and back-propagating through all $K$ steps is expensive and a known route to reward hacking, the estimator needs a critic that scores a noisy intermediate state, trained jointly with the policy from verifier outcomes and back-propagated through the denoiser at one random flow step. With the scalar critic of prior reward back-propagation work this already reaches the highest peaks we observed, but only at a learning rate at which it is fragile: some seeds never learn and others collapse.

The third level is our contribution, and it starts from where credit should be indexed. In AR post-training every quantity the algorithm manipulates, the action, the advantage, the value, is indexed by the token position alone, because the token position is the RL time step; in a flow model a generation is a grid of states over flow steps and token positions, and both the ported GRPO and the scalar critic keep a single axis, the flow step, discarding the token axis on which the semantics live. We therefore build a \hl{two-dimensional (2D) critic} (Figure~\ref{fig:mechanism}), a learned value over both axes: the \emph{flow step}, i.e.\ \emph{when} along the trajectory a noisy state is evaluated, and the \emph{token position}, i.e.\ \emph{where} along the sequence the outcome is decided. It outputs one value per position of any intermediate state and pools them into the value the policy ascends; the scalar critic is the special case that pools before it evaluates, and the 2D form is what makes joint training stable at a single learning rate across tasks and seeds. We call the resulting actor--critic recipe \jpc{} (joint pathwise critic).

Around these three levels we assemble a complete post-training study: the estimator analysis, the three RL recipes under one evaluation, an amortisation stage (\S\ref{sec:loop}) that distils the RL policy's verified samples back into the model so that samplers far cheaper than the one used for RL inherit the improvement, and two external controls on the same data, an AR model with the standard recipe and the masked-diffusion recipe d1 \citep{zhao2025d1}. We evaluate on four reasoning tasks with program verifiers (Countdown, Sudoku, Knights-and-Knaves, and a Game24 split that tests generalisation to unseen number sets), three seeds per RL recipe.

\paragraph{Contributions.}
(i) A structural argument, with matching measurements, for why score-function RL delivers a weak and noisy signal on continuous flow language models and why pathwise gradients are their native estimator (\S\ref{sec:estimators}).
(ii) The \hl{2D critic}, a value over flow steps and token positions that turns the pathwise gradient into a practical actor--critic method, \jpc{}, whose ingredients (per-position values pooled to a state value, joint training, the deployment sampler, one random step per update, early stopping on a held-out set) we ablate individually (\S\ref{sec:method}, \S\ref{sec:ablations}).
(iii) A three-level comparison on four tasks with three seeds each (\S\ref{sec:experiments}): GRPO improves every task but is seed-sensitive and gains least where there is most to learn; the scalar-critic pathwise recipe reaches the highest peaks but is fragile; the 2D critic replicates on every task and seed, matches GRPO where the headroom is small and leads clearly where it is large.
(iv) An amortisation stage that carries the RL gains to cheap samplers, with controls that attribute the gains to the RL policy, and external baselines (an AR model with token-level PPO; d1 on LLaDA-8B) on the same data (\S\ref{sec:loop}, \S\ref{sec:experiments}, Appendix~\ref{app:d1}).
